\begin{figure}[H]
\centering
\begin{tikzpicture}[>=Latex,scale=0.92]
    % Panel izquierdo: tipos de líneas de universo
    \begin{scope}[xshift=-5.0cm]
        \node[font=\bfseries] at (0,4.35) {(a) Tipos de líneas de universo};
        \draw[->,black,very thick] (-3.7,-3.35) -- (3.85,-3.35)
            node[right] {$x$};
        \draw[->,black,very thick] (-3.35,-3.7) -- (-3.35,3.9)
            node[above] {$ct$};

        \draw[black,very thick] (-2.45,-3.25) -- (-2.45,3.35);
        \node[black,font=\small,align=center] at (-2.45,3.72)
            {reposo};

        \draw[blue!75!black,very thick] (-1.45,-3.25) -- (0.15,3.35);
        \node[blue!75!black,font=\small,align=center] at (0.15,3.72)
            {MRU};

        \draw[red!75!black,very thick]
            (0.25,-3.25) .. controls (0.2,-1.3) and (1.05,0.45) .. (2.65,3.35);
        \node[red!75!black,font=\small,align=center] at (2.45,3.72)
            {aceleración};

        \draw[orange!85!black,very thick,dashed] (-1.2,-3.25) -- (3.65,1.6);
        \node[orange!85!black,font=\small,rotate=45] at (2.65,0.95)
            {luz};

        \node[gray!65!black,font=\scriptsize,align=center] at (0,-4.1)
            {una partícula masiva tiene una trayectoria tipo tiempo};
    \end{scope}

    % Panel derecho: conos locales sobre una trayectoria
    \begin{scope}[xshift=4.9cm]
        \node[font=\bfseries] at (0,4.35) {(b) Conos de luz locales};
        \draw[->,black,very thick] (-3.7,-3.35) -- (3.85,-3.35)
            node[right] {$x$};
        \draw[->,black,very thick] (-3.35,-3.7) -- (-3.35,3.9)
            node[above] {$ct$};

        \coordinate (Eone) at (-1.45,-1.75);
        \coordinate (Etwo) at (0.05,0.55);
        \coordinate (Ethree) at (1.25,2.35);

        \foreach \event in {Eone,Etwo,Ethree}{
            \fill[blue!8]
                (\event) -- ($(\event)+(-0.72,0.72)$)
                arc[start angle=180,end angle=360,
                    x radius=0.72,y radius=0.18] -- cycle;
            \fill[blue!8]
                (\event) -- ($(\event)+(-0.72,-0.72)$)
                arc[start angle=180,end angle=0,
                    x radius=0.72,y radius=0.18] -- cycle;

            \draw[blue!65!black,thick]
                (\event) -- ($(\event)+(-0.72,0.72)$);
            \draw[blue!65!black,thick]
                (\event) -- ($(\event)+(0.72,0.72)$);
            \draw[blue!65!black,thick]
                (\event) -- ($(\event)+(-0.72,-0.72)$);
            \draw[blue!65!black,thick]
                (\event) -- ($(\event)+(0.72,-0.72)$);
            \draw[blue!65!black,thick]
                ($(\event)+(0,0.72)$) ellipse (0.72 and 0.18);
            \draw[blue!65!black,thick]
                ($(\event)+(0,-0.72)$) ellipse (0.72 and 0.18);
        }

        \draw[magenta!60!black,very thick]
            plot[smooth,tension=0.65] coordinates {
                (-1.7,-3.1)
                (-1.45,-1.75)
                (0.05,0.55)
                (1.25,2.35)
                (1.55,3.35)
            };

        \foreach \event/\name in {Eone/E_1,Etwo/E_2,Ethree/E_3}{
            \filldraw[black] (\event) circle (2.2pt);
            \node[anchor=west,font=\small] at ($(\event)+(0.13,-0.18)$)
                {$\name$};
        }

        \node[magenta!60!black,font=\small,align=center] at (2.35,-2.7)
            {línea de universo\\de una partícula masiva};
        \draw[->,magenta!60!black,thin] (1.65,-2.45) -- (-0.85,-1.0);
        \node[blue!65!black,font=\small,align=center] at (2.45,3.45)
            {cada evento tiene\\su cono local};
    \end{scope}
\end{tikzpicture}
\caption{Líneas de universo y causalidad local. (a) Trayectorias de un objeto en reposo, con movimiento rectilíneo uniforme, con aceleración y de un pulso de luz. (b) Una línea de universo masiva atraviesa una sucesión de eventos, cada uno con su propio cono de luz; la trayectoria permanece siempre dentro de los conos locales.}
\label{fig:lineas_universo}
\end{figure}